\section{Method}

\subsection{Problem Formulation}
We control a fleet of $N$ vehicles on a fixed SUMO road network while a background
population of uncontrolled vehicles shares the roads. Each controlled vehicle $i$ is an
agent that, at a set of route-decision points along its trip, chooses a complete route to
its destination. SUMO handles all low-level dynamics---car-following, lane changing,
junctions, and signals---so the policy only sets high-level route intent. The fleet
objective is to minimize average travel time,
\begin{equation}
  J(\pi) = \frac{1}{N}\sum_{i=1}^{N} T_i,
\end{equation}
where $T_i$ is the realized travel time of vehicle $i$ under the shared policy $\pi$.
This is a cooperative, partially observed problem: each agent sees only a local view of
the network but its route choice affects the congestion that other agents experience.
We treat it as a decentralized POMDP and solve it with CTDE MAPPO, so that a fleet-wide
signal can guide learning while deployment stays per-vehicle.

The tension the policy must resolve is precisely the price-of-anarchy gap of
Section~II: the route that minimizes $T_i$ for one vehicle need not minimize $J$ for the
fleet. A useful policy learns when to keep a vehicle on its shortest feasible path
(most of the time) and when to route it away from a forming bottleneck for the group's
benefit.

\begin{figure}[!t]
\centering
\resizebox{\columnwidth}{!}{%
\begin{tikzpicture}[
    node distance=0.72cm,
    box/.style={rectangle, draw, rounded corners, align=center, minimum width=3.0cm, minimum height=0.62cm, font=\footnotesize},
    arrow/.style={->, thick}
]
    \node[box] (sumo) {SUMO microsimulation\\(local + fleet state)};
    \node[box, below=of sumo] (cand) {Route-candidate generator\\(up to 4 routes, guardrail filters)};
    \node[box, below=of cand] (actor) {MAPPO actor\\score \& mask candidates};
    \node[box, below=of actor] (throttle) {Coordination guardrails\\(detour guard, reservations)};
    \node[box, below=of throttle] (exec) {Route commit\\+ SUMO execution};
    \node[box, below=of exec] (learn) {Reward + centralized critic\\GAE / PPO update \emph{(training only)}};
    \draw[arrow] (sumo) -- (cand);
    \draw[arrow] (cand) -- (actor);
    \draw[arrow] (actor) -- (throttle);
    \draw[arrow] (throttle) -- (exec);
    \draw[arrow] (exec) -- (learn);
    \draw[arrow] (learn.east) to[out=0,in=0,looseness=1.7] node[right, font=\scriptsize]{next step} (sumo.east);
\end{tikzpicture}%
}
\caption{The STR-SUMO learning and deployment pipeline. Every stage speaks a common
travel-time / density vocabulary. The dashed portion (reward and critic) runs only during
training; at deployment the actor and guardrails alone drive each vehicle.}
\label{fig:pipeline}
\end{figure}

\subsection{Route-Candidate Action Space}
Rather than asking the policy to pick a turn (``left/right'') at each junction, a
route-candidate generator proposes up to $K{=}4$ \emph{complete} routes from the vehicle's
current edge to its destination. Index~0 is always the congestion-aware shortest route;
indices 1--3 are diverse detours produced by density-penalized Dijkstra searches. The
actor's job is to choose one of these routes, so a ``selfless detour'' simply means
choosing an index other than~0. Each candidate carries an 11-dimensional feature
``scorecard'' describing, among other things, its average and worst-case congestion, its
first-edge congestion, the extra seconds it costs relative to the shortest route (the
ETA sacrifice), and how much less congested it is than the shortest route (the relief).
Because the reward, the actor, and the guardrails all read the same scorecard, a
trade-off assessed by one component is understood identically by the others.

\subsection{CTDE Actor and Critic}
Figure~\ref{fig:ctde} shows the network. The \textbf{actor} is a shared per-candidate
scorer: a small multilayer perceptron (two hidden layers of 256 and 128 units, ReLU)
takes the vehicle's local observation concatenated with one candidate's features and
emits a scalar score; the same sub-network scores all $K$ candidates, and invalid
candidates are masked to $-10^{9}$ before a softmax over scores. This weight-sharing lets
the actor generalize across a variable, unordered candidate set and keeps deployment
per-vehicle. The local observation encodes the current and destination edges, feasibility
masks, lane and commit-window state, route-difficulty estimates, local density summaries,
and the per-candidate features.

The \textbf{critic} is used only during training. It concatenates the local observation
with an 18-dimensional \emph{central} observation summarizing the whole fleet's state
(aggregate densities, completion, congestion pressure) and predicts a scalar value. This
is the CTDE split~\cite{lowe2017maddpg}: the critic ``peeks'' at global information to
give lower-variance value estimates, while the actor never does, so the trained policy
runs decentralized. The critic changes learning stability, not the objective; the
cooperative incentive must come from the reward.

\begin{figure}[!t]
\centering
\resizebox{\columnwidth}{!}{%
\begin{tikzpicture}[
    font=\scriptsize,
    box/.style={rectangle, draw, rounded corners, align=center, minimum height=0.55cm, inner sep=3pt},
    sm/.style={rectangle, draw, align=center, minimum height=0.5cm, inner sep=2pt, fill=blue!5},
    arrow/.style={->, thick}
]
    \node[box, fill=orange!8] (obs) {local obs $o_i$};
    \node[sm, right=0.5cm of obs] (c0) {cand.\ $k$\\features};
    \node[box, below=0.55cm of obs, minimum width=3.4cm] (scorer) {shared candidate scorer (MLP)};
    \node[box, below=0.5cm of scorer, minimum width=3.4cm] (mask) {mask invalid $\rightarrow$ softmax over $K$};
    \node[box, below=0.5cm of mask, fill=orange!8, minimum width=3.4cm] (act) {sampled / argmax route index};
    \draw[arrow] (obs.south) -- ([xshift=-1.0cm]scorer.north);
    \draw[arrow] (c0.south) -- ([xshift=1.0cm]scorer.north);
    \draw[arrow] (scorer) -- (mask);
    \draw[arrow] (mask) -- (act);

    \node[box, fill=blue!8, right=2.4cm of obs] (cobs) {central obs $s$ (18-d)};
    \node[box, below=1.1cm of cobs, minimum width=2.6cm] (critic) {critic MLP};
    \node[box, below=0.5cm of critic, minimum width=2.6cm, fill=blue!8] (val) {value $V(o_i,s)$};
    \draw[arrow] (cobs.south) -- (critic.north);
    \draw[arrow] (obs.east) |- ([yshift=0.15cm]critic.west);
    \draw[arrow] (critic) -- (val);
    \node[align=center, below=0.28cm of act, font=\scriptsize\itshape] {actor (train + deploy)};
    \node[align=center, below=0.28cm of val, font=\scriptsize\itshape] {critic (train only)};
\end{tikzpicture}%
}
\caption{CTDE MAPPO architecture. The actor scores each candidate route with a shared
sub-network from the vehicle's local view; the centralized critic additionally reads an
18-dimensional fleet summary and is discarded at deployment.}
\label{fig:ctde}
\end{figure}

\subsection{Learning Algorithm}
The policy is trained with MAPPO~\cite{yu2022mappo}: a clipped PPO surrogate for the
actor~\cite{schulman2017ppo} and a squared-error value loss for the critic, with an
entropy bonus for exploration and separate Adam optimizers~\cite{kingma2015adam} for the
two networks. Advantages are computed per vehicle trajectory with
GAE~\cite{schulman2016gae} ($\gamma{=}0.995$, $\lambda{=}0.95$); the high discount is
needed because trips span roughly two thousand simulation steps and a route choice must
be credited for congestion it causes far downstream. Two practical components proved
important on this task. First, because episode returns swing by more than an order of
magnitude with the demand seed, the critic regresses on \emph{normalized} value targets
(running mean/variance), while advantages remain in real reward units; this keeps the
value loss well-conditioned. Second, entropy is annealed from 0.15 to 0.05 over training
so the policy explores detours early and settles later, and PPO updates are stopped early
when the approximate KL divergence exceeds 0.015. Key hyperparameters are listed in
Table~\ref{tab:hyper}.

\begin{table}[!t]
\caption{MAPPO Training Configuration}
\label{tab:hyper}
\centering
\footnotesize
\begin{tabular}{lll}
\toprule
Group & Parameter & Value \\
\midrule
Actor/critic & hidden layers & $[256,128]$, ReLU \\
             & candidate routes $K$ & 4 \\
             & central obs dim & 18 \\
Optimization & actor / critic LR & $3\times10^{-4}$ / $1\times10^{-3}$ \\
             & $\gamma$ / GAE $\lambda$ & 0.995 / 0.95 \\
             & PPO clip $\epsilon$ & 0.05 \\
             & epochs / minibatch & 4 / 512 \\
             & target KL & 0.015 \\
             & value-target norm. & on \\
Exploration  & entropy coef.\ (start$\to$end) & $0.15 \to 0.05$ \\
Objective    & team-reward $\alpha$ & 1.0 \\
             & team-reward mode & difference \\
\bottomrule
\end{tabular}
\end{table}

\subsection{Reward: Cooperation in Travel-Time Units}
Every term of the reward is measured in seconds so trade-offs need no arbitrary
conversion factor. Each simulation step charges a small base time cost (0.07 per step),
plus deliberately smaller costs for occupying above-average-density roads and for the
congestion a vehicle imposes on others; reaching the destination pays $+50$ and a
teleport (SUMO's last resort for a stuck vehicle) costs $-40$. The cooperative term is
the key addition. Let $\rho_i$ be vehicle $i$'s normalized slowness (0 at free flow, 1 at
standstill) and $\bar\rho$ the live-fleet mean. With weight $\alpha\in[0,1]$ and scale
$\kappa$, agent $i$ pays over an elapsed interval $\Delta t$
\begin{equation}
  c^{\text{team}}_i = \alpha\,\kappa\,(\rho_i - \bar\rho)\,\Delta t .
\end{equation}
This is a leave-one-out difference reward~\cite{wolpert2001collectives}: because the
fleet-mean of $(\rho_i-\bar\rho)$ is zero, the exogenous, demand-driven congestion level
cancels, and each agent is credited only for how its own route choice fares relative to
the fleet---a vehicle that routed onto a free-flowing alternative is rewarded, one that
piled into a jam pays. We use $\alpha{=}1$ (fleet delay counts as much as personal delay)
with the difference mode throughout. Setting $\alpha{=}0$ recovers a purely selfish
objective and serves as a conceptual control.

\subsection{Guardrails and the Saturation Veto}
Hand-written guardrails keep the reward focused on travel time by removing obviously bad
options rather than teaching basic safety through penalties. Pre-decision filters prune
candidate routes that would loop, re-enter a dead end, or drive straight into a known jam;
a ``pending'' state machine distinguishes healthy queueing (e.g.\ at a red light) from a
genuine stall so the reward is not charged for normal waiting. Two coordination guardrails
address a specific failure in heavy congestion, where independently chosen detours can pile
onto the same alternative and make it worse:
\begin{itemize}
  \item The \textbf{detour guard} (deployment only): after the actor picks a detour, if the
  alternative is itself near capacity and the network is congested, the choice is reverted
  to the shortest route. It judges the alternative by its \emph{absolute} crowdedness, not
  by the relief the policy reacted to, because when everything is busy a road can look like
  relief yet still have no room left. It runs only at deployment, since overriding an
  executed route during training would break the on-policy assumption.
  \item The \textbf{route reservation} (training and deployment): a committed route books
  its leading edges in a decaying field, so the candidate generator scores later deciders
  against effective (live~+~reserved) density and a simultaneous rush becomes an orderly
  one-at-a-time spread. It affects only what candidates look like, never the reward or the
  true measurements.
\end{itemize}
The detour guard is the subject of our ablation in Section~V-D: it decides whether the
policy's learned detours are executed or suppressed, and it turns out to trade peak fleet
gains for robustness.

\section{Experimental Setup}

\subsection{Network and Demand}
Experiments use a corner of the New York City street grid, cut from
OpenStreetMap~\cite{osm} and converted to a SUMO network with SUMO's OSM import
tooling~\cite{lopez2018sumo}: 132 drivable edges, 360 lanes, 90 junctions (53
signalized), and roughly 27{,}000 lane-meters. To create genuine
congestion---a precondition for selfless routing to have anything to gain---we use a
shared-destination demand pattern (ranged origins routed to one common destination edge)
with 450 controlled and 150 background vehicles released at a 0.5\,s interval. This
concentrates flow into a corridor and produces sustained congestion (frequent teleports,
fleet travel times from a few hundred to over 1500\,s across seeds) without permanent
gridlock. A dispersed origin--destination pattern, by contrast, leaves the grid at a few
percent of capacity, where routing cannot matter.

\subsection{Protocol, Baselines, and Metrics}
The policy is trained for 100 episodes. Every 25 episodes it is evaluated frozen (no
learning) on 11 held-out seeds (6000--6010) against a Dijkstra baseline generated from the
identical demand, and the best checkpoint is retained. We then \emph{deploy} the best
checkpoint on a separate block of 20 held-out scenarios (seeds 4010--4029) and compare it
head-to-head with Dijkstra on the same generated vehicles. We report mean fleet travel
time, median (p50) and p90 travel time, tail-completion gap, the p95/p50 ratio, deadline
misses, completion rate, and per-scenario win rate. Unless noted, the MAPPO policy is
deployed \emph{greedily} (argmax over candidate scores), which is deterministic and, as
Section~V-E shows, matches stochastic sampling once the detour guard is active. Each
controller is run in a fresh simulation with independent network state, so no traffic
leaks between runs.
